\documentclass[a4paper, serif, 10pt]{moderncv}

%% moderncv
% color and style
\moderncvcolor{black}
\moderncvstyle{banking}

% renew moderncv command to adapt for CJK colon
\renewcommand*{\cvitem}[3][.25em]{%
  \ifstrempty{#2}{}{\hintstyle{#2}：}{#3}%
  \par\addvspace{#1}}

\renewcommand*{\cvitemwithcomment}[4][.25em]{%
  \savebox{\cvitemwithcommentmainbox}{\ifstrempty{#2}{}{\hintstyle{#2}：}#3}%
  \setlength{\cvitemwithcommentmainlength}{\widthof{\usebox{\cvitemwithcommentmainbox}}}%
  \setlength{\cvitemwithcommentcommentlength}{\maincolumnwidth-\separatorcolumnwidth-\cvitemwithcommentmainlength}%
  \begin{minipage}[t]{\cvitemwithcommentmainlength}\usebox{\cvitemwithcommentmainbox}\end{minipage}%
  \hfill% fill of \separatorcolumnwidth
  \begin{minipage}[t]{\cvitemwithcommentcommentlength}\raggedleft\small\itshape#4\end{minipage}%
  \par\addvspace{#1}}

%% page layout/margins
\usepackage[top=2.5cm, bottom=2.5cm, left=1.5cm, right=1.5cm]{geometry}

\nopagenumbers{}

%% Babel config for English language
\usepackage[english]{babel}

%% fontspec
\usepackage{fontspec}

\IfFontExistsTF{Linux Libertine}{
  \setmainfont[Ligatures={TeX, Common}, Numbers=Lining, ItalicFont=Linux Libertine]{Linux Libertine}
}{}
\IfFontExistsTF{Linux Libertine O}{
  \setmainfont[Ligatures={TeX, Common}, Numbers=Lining, ItalicFont=Linux Libertine O]{Linux Libertine O}
}{}

%% CTeX
% CJK support, used to show CJK characters in the resume
%
% - heading=false: disable ctex heading
% - punct=kaiming: use kaiming punctuations styles for CJK
% - scheme=plain: use plain scheme, do not override \normalsize font size
% - space=auto: space settings for CJK characters
%
% CJK fonts are resolved per locale language, e.g. Noto Serif/Sans CJK JP
% for Japanese resumes, with fallback to the next region if not installed
%
% ref:
% - http://ctan.mirrorcatalogs.com/language/chinese/ctex/ctex.pdf
\usepackage[UTF8, heading=false, punct=kaiming, scheme=plain, space=auto]{ctex}

\IfFontExistsTF{Noto Serif CJK TC}{
  \setCJKmainfont{Noto Serif CJK TC}
}{}
\IfFontExistsTF{Noto Sans CJK TC}{
  \setCJKsansfont{Noto Sans CJK TC}
}{}

%% line spacing
\usepackage{setspace}
\setstretch{1.125}

%% URL styling - use normal text instead of monospace
\urlstyle{same}

%% Patch \httplink and \httpslink to handle full URLs
% The original moderncv macros always prepend http:// or https:// to URLs.
% This causes issues when the URL already has a protocol (e.g., https://example.com)
% resulting in malformed URLs like https://https://example.com.
% This patch checks if the URL already contains "://" and uses it directly if so.
% Note: We use \str_set:Nx (x-expansion) to expand macros like \@homepage first.
\ExplSyntaxOn
\str_new:N \g_yamlresume_url_str

\RenewDocumentCommand{\httpslink}{O{}m}{
  \str_set:Nx \g_yamlresume_url_str {#2}
  \str_if_in:NnTF \g_yamlresume_url_str {://}
    { \url{#2} }
    { \url{https://#2} }
}

\RenewDocumentCommand{\httplink}{O{}m}{
  \str_set:Nx \g_yamlresume_url_str {#2}
  \str_if_in:NnTF \g_yamlresume_url_str {://}
    { \url{#2} }
    { \url{http://#2} }
}
\ExplSyntaxOff

\name{林冠宇}{}
\title{資深軟體測試工程師}
\phone[mobile]{+886 912 345 678}
\email{kuanyu.lin@example.com}
\homepage{https://kuanyulin.dev}

\address{中國臺灣，台北市，台北市，台北市信義區松仁路 100 號，110}{}{}

\extrainfo{{\small \faLinkedin}\ \href{https://www.linkedin.com/in/kuanyulin}{@kuanyulin} {} {} {} • {} {} {} 
{\small \faGithub}\ \href{https://github.com/kuanyulin}{@kuanyulin}}

\begin{document}

\maketitle

\section{簡介}

\cvline{}{\begin{itemize}
\item 擁有 8 年軟體測試經驗，專精於自動化測試框架設計與導入
\item 熟悉 Web、API 與行動裝置應用程式測試，具備紮實的程式開發能力
\item 擅長建立測試流程與品質指標，推動 CI/CD 與測試左移
\item 具備跨團隊溝通與專案管理能力，能有效提升產品品質與交付效率
\end{itemize}}
\section{教育背景}

\cventry{2012年9月–2016年6月}
        {學士，資訊工程學系，成績：3.9/4.0}
        {國立臺灣大學}
        {\url{https://www.ntu.edu.tw}}
        {}
        {\begin{itemize}
\item 主修資訊工程，奠定紮實的程式設計與系統分析基礎
\item 專題研究聚焦於軟體測試與品質保證，獲得系上專題競賽優等
\item 擔任系學會資訊組組長，負責系網維護與活動系統開發
\end{itemize}
\textbf{課程}：資料結構、演算法、作業系統、資料庫系統、軟體工程、計算機網路、程式設計、離散數學}
\section{工作經歷}

\cventry{2020年1月至今}
        {資深品質保證工程師}
        {趨勢科技}
        {\url{https://www.trendmicro.com}}
        {}
        {\begin{itemize}
\item 負責雲端安全產品（CWPP）的整體測試策略規劃與執行，涵蓋功能、效能、安全性與相容性測試
\item 主導建置以 Pytest 與 Selenium 為基礎的端到端自動化測試框架，將回歸測試時間從 3 天縮短至 4 小時
\item 導入測試左移流程，與開發團隊協作於 CI/CD 管道中執行單元測試與整合測試，提升版本交付品質
\item 建立品質指標與測試覆蓋率報告，定期向產品與工程團隊彙報，推動持續改善
\item 指導 3 位初階測試工程師，協助其熟悉自動化測試與缺陷管理流程
\end{itemize}
\textbf{關鍵字}：自動化測試、雲端安全、Pytest、Selenium、CI/CD、測試策略}

\cventry{2016年7月–2019年12月}
        {品質保證工程師}
        {聯發科技}
        {\url{https://www.mediatek.com}}
        {}
        {\begin{itemize}
\item 負責行動晶片平台之相容性與效能測試，確保產品符合業界標準與客戶需求
\item 設計並執行測試計畫，涵蓋 Android 相容性測試（CTS）、效能基準測試與功耗測試
\item 開發自動化測試腳本（Java、Python），整合 Jenkins 以實現每日建置與測試
\item 分析測試結果並追蹤缺陷，與硬體、軟體及客戶團隊協作解決問題
\item 撰寫測試報告與技術文件，協助提升團隊測試效率與知識分享
\end{itemize}
\textbf{關鍵字}：嵌入式測試、效能測試、Android CTS、Java、Jenkins、缺陷追蹤}
\section{語言}

\cvline{中文}{母語或雙語水平 \hfill \textbf{關鍵字}：母語}
\cvline{英語}{完全專業水平 \hfill \textbf{關鍵字}：TOEIC 950、技術文件撰寫、跨國團隊溝通}
\section{專業技能}

\cvline{測試自動化}{專家 \hfill \textbf{關鍵字}：Selenium、Cypress、Appium、Pytest、Python}
\cvline{效能與負載測試}{高級 \hfill \textbf{關鍵字}：JMeter、Locust、Gatling、k6}
\cvline{API 測試}{高級 \hfill \textbf{關鍵字}：Postman、REST Assured、Swagger、GraphQL}
\cvline{測試管理與流程}{高級 \hfill \textbf{關鍵字}：JIRA、TestRail、Agile、Scrum、CI/CD}
\cvline{程式與資料庫}{中級 \hfill \textbf{關鍵字}：Python、Java、SQL、Git、Docker}
\section{榮譽}

\cventry{2023年12月}
        {趨勢科技}
        {年度卓越員工獎}
        {}
        {}
        {表揚在雲端安全產品測試自動化與品質流程改善上的傑出貢獻，成功將回歸測試時間縮短 90\%。}

\cventry{2021年11月}
        {台灣軟體測試社群}
        {黑客松最佳測試創新獎}
        {}
        {}
        {以「智慧型測試案例生成器」為題，運用機器學習分析使用者行為並自動產生高覆蓋率測試案例，獲得評審青睞。}
\section{證書}

\cventry{2021年3月}
        {ISTQB}
        {ISTQB 進階測試經理認證}
        {\url{https://www.istqb.org/certifications/advanced-level-test-manager}}
        {}
        {}

\cventry{2022年8月}
        {AWS}
        {AWS 雲端從業人員認證}
        {\url{https://aws.amazon.com/certification/certified-cloud-practitioner/}}
        {}
        {}

\cventry{2020年6月}
        {Scrum Alliance}
        {認證 ScrumMaster}
        {\url{https://www.scrumalliance.org/get-certified/scrum-master-track/certified-scrummaster}}
        {}
        {}
\section{出版刊物}

\cventry{2022年5月}
        {建構高覆蓋率自動化測試框架的實務指南}
        {軟體測試技術月刊}
        {\url{https://www.softwaretesting.tw/articles/automation-framework-guide}}
        {}
        {\begin{itemize}
\item 探討如何從零開始建構可維護、可擴展的自動化測試框架
\item 分享實際案例與程式碼範例，涵蓋 Page Object Model 與資料驅動測試
\item 提供測試覆蓋率與品質指標的衡量方法
\end{itemize}}
\section{項目}

\cventry{2021年1月–2022年12月}
        {為雲端安全產品設計與實作端到端自動化測試框架}
        {雲端安全產品自動化測試框架}
        {\url{https://github.com/kuanyulin/cloud-security-test-framework}}
        {}
        {\begin{itemize}
\item 採用 Pytest 與 Selenium 建構模組化測試框架，支援多環境與多瀏覽器執行
\item 整合 Docker 與 Jenkins，實現每日自動化回歸測試與報告產出
\item 測試覆蓋率提升至 85\%，缺陷漏檢率降低 40\%
\end{itemize}
\textbf{關鍵字}：自動化測試、Pytest、Selenium、Docker、Jenkins}

\cventry{2019年3月–2019年11月}
        {開發輕量級行動裝置效能測試工具，用於監控 CPU、記憶體與電量消耗}
        {行動裝置效能測試工具}
        {\url{https://github.com/kuanyulin/mobile-perf-tool}}
        {}
        {\begin{itemize}
\item 使用 Python 與 ADB 開發效能資料收集與視覺化工具
\item 支援多款 Android 裝置，協助團隊快速定位效能瓶頸
\item 工具導入後，效能問題平均修復時間縮短 30\%
\end{itemize}
\textbf{關鍵字}：效能測試、Python、ADB、Android、資料視覺化}
\section{興趣愛好}

\cvline{登山}{百岳、健行、露營}
\cvline{攝影}{風景攝影、街拍}
\cvline{閱讀}{技術書籍、科幻小說}
\section{志願服務}

\cventry{2021年1月–2023年12月}
        {志工講者}
        {台灣軟體測試社群}
        {\url{https://www.softwaretesting.tw}}
        {}
        {\begin{itemize}
\item 定期於社群 meetup 分享自動化測試實務經驗與最佳實踐
\item 協助籌辦年度軟體測試研討會，吸引超過 300 位參與者
\item 參與開源測試工具推廣，翻譯技術文件並貢獻測試案例
\end{itemize}}
    
\end{document}